\section{符号说明}
为避免四个问题分别定义不同符号，本文按照“公共物理量—基础决策变量—后续新增业务变量”的层级统一记号。

\subsection{下标与集合}

\begin{table}[H]
\centering
\zihao{5}
\begin{tabular}{lll}
\toprule
符号 & 含义 & 来源 \\
\midrule
\(d\) & 日期索引 & 全年滚动计算需要区分不同日期 \\
\(t\) & 日内第 \(t\) 个 10 min 时段，\(t=1,\ldots,144\) & 由一天按 10 min 离散得到 \\
\(\tau\) & 滚动视野中的未来日期 & 7 日滚动优化需要区分当前日与未来日 \\
\(s\) & 日内调整阶段，\(s\in\{0,6,12,18\}\) & 问题三、四的预报发布时间 \\
\(\omega\) & SAA 情景编号，\(\omega=1,\ldots,K\) & 由历史预测误差抽样产生 \\
\(K\) & SAA 情景数，正式取 \(K=4\) & 随机规划的样本数量 \\
\(R\) & 滚动视野长度，正式取 \(R=7\) 天 & 用于体现日末储电量的未来价值 \\
\bottomrule
\end{tabular}
\caption{下标与集合符号说明}
\label{tab:1}
\end{table}

\subsection{已知参数与输入数据}

\begin{table}[H]
\centering
\zihao{5}
\begin{tabular}{llr}
\toprule
符号 & 含义 & 单位 \\
\midrule
\(L_{d,t}\) & 小区负荷功率 & kW \\
\(PV_{d,t}\) & 光伏发电功率 & kW \\
\(\pi_t\) & 固定日内电价（问题一至三） & 元/kWh \\
\(\pi^{act}_{d,t}\) & 问题四附件 4 已实现真实电价 & 元/kWh \\
\(\hat\pi_{d,t}\) & 计划阶段的电价中心预测 & 元/kWh \\
\(\pi^{(\omega,s)}_{d,t}\) & 阶段 \(s\) 下第 \(\omega\) 个价格情景 & 元/kWh \\
\(E_{\min}\) & 储能运行电量下限，1200 & kWh \\
\(E_{\max}\) & 储能运行电量上限，10800 & kWh \\
\(P_{\max}\) & 最大充放电功率，5000 & kW \\
\(\eta_c,\eta_d\) & 充电效率、放电效率，均为 0.90 & — \\
\(\Delta t\) & 单个调度时段长度，\(1/6\) & h \\
\bottomrule
\end{tabular}
\caption{已知参数与输入数据（含单位）}
\label{tab:2}
\end{table}

\subsection{由输入数据直接构造的辅助量}

光伏优先满足当前负荷，因此首先从 \(L\) 与 \(PV\) 构造三类不需要优化的辅助量：

\[
PV^L_{d,t}=\min(PV_{d,t},L_{d,t}),
\]

表示光伏直接供负荷功率；

\[
\overline L_{d,t}
=
\max(L_{d,t}-PV_{d,t},0),
\]

表示光伏供电后仍存在的负荷缺口；

\[
\overline{PV}_{d,t}
=
\max(PV_{d,t}-L_{d,t},0),
\]

表示满足负荷后的富余光伏。

这三个量来源于“光伏优先供负荷”的调度规则，因此由输入数据预先计算，而不是交给优化器自由决定。

\subsection{四问共用的基础物理决策变量}

\begin{table}[H]
\centering
\zihao{5}
\begin{tabular}{lll}
\toprule
变量 & 含义 & 为什么需要该变量 \\
\midrule
\(G^L\) & 正常外网购电中直接供负荷的功率 & 区分外网购电用于负荷还是用于储能 \\
\(G^{ch}\) & 正常外网购电中用于储能充电的功率 & 刻画低价购电充电行为 \\
\(PV^{ch}\) & 富余光伏用于储能充电的功率 & 刻画光伏余电消纳 \\
\(C\) & 储能充电功率 & 储能控制变量 \\
\(D\) & 储能放电功率 & 储能控制变量 \\
\(E\) & 时段结束后的储电量 & 描述储能状态并进行跨时段递推 \\
\(V\) & 弃光功率 & 不允许售电时必须给富余光伏设置去向 \\
\(u\) & 储能充放电状态二元变量 & 强制同一时段不同时充放电 \\
\bottomrule
\end{tabular}
\caption{由输入数据直接构造的辅助量}
\label{tab:3}
\end{table}

正常外网总购电功率不再单独设置自由变量，而由

\[
G=G^L+G^{ch}
\]

派生得到。这样能够直接追踪每一单位外网电的物理去向。

\subsection{问题二新增变量}

问题二新增“计划时刻”和“实际执行时刻”的信息差，因此在基础物理变量之外增加：

\begin{table}[H]
\centering
\zihao{5}
\begin{tabular}{lll}
\toprule
变量 & 含义 & 变量来源 \\
\midrule
\(G^{plan}_{d,t}\) & 每天 0:00 提交的正常计划购电功率 & 题目要求事前制定全天计划 \\
\(H_{d,t}\) & 紧急购电功率 & 实际供电不足时按 5 倍电价补购 \\
\(W_{d,t}\) & 已购未用正常电功率 & 计划购电已锁定但实际无法利用时需要守恒去向 \\
\bottomrule
\end{tabular}
\caption{四问共用的基础物理决策变量}
\label{tab:4}
\end{table}

因此当天计划与执行之间满足

\[
G^L+G^{ch}+W=G^{plan}.
\]

\subsection{问题三新增变量}

问题三允许在 6:00、12:00、18:00 调整 0:00 原始计划，因此新增：

\begin{table}[H]
\centering
\zihao{5}
\begin{tabular}{lll}
\toprule
变量 & 含义 & 变量来源 \\
\midrule
\(P_{d,t}\) & 0:00 原始计划购电功率 & 当天最初提交的计划 \\
\(A^s_{d,t}\) & 阶段 \(s\) 调整后的购电功率 & 新光伏预报到达后的再决策 \\
\(B_{d,t}\) & 全天最终生效计划 & 将四阶段实际生效区间拼接得到 \\
\(\Delta^+_{d,t}\) & 相对原计划的最终调增量 & 1.5 倍电价的调增结算 \\
\(\Delta^-_{d,t}\) & 相对原计划的最终调减量 & 0.5 倍电价的违约结算 \\
\bottomrule
\end{tabular}
\caption{问题二新增的随机规划变量}
\label{tab:5}
\end{table}

其中

\[
B-P=\Delta^+-\Delta^-.
\]

问题三固定电价为正，最优解下 \(\Delta^+,\Delta^-\) 不会同时为正，因此无需额外互斥变量。

\subsection{问题四新增变量}

问题四将未来价格本身作为随机量，因此新增：

\begin{table}[H]
\centering
\zihao{5}
\begin{tabular}{lll}
\toprule
变量 & 含义 & 变量来源 \\
\midrule
\(\hat\pi\) & 电价中心预测 & 未来真实价格不可提前获得 \\
\(e^\pi\) & 历史价格预测残差 & 用于生成价格 SAA 情景 \\
\(\pi^{(\omega,s)}\) & 价格随机情景 & 由中心预测与历史残差叠加 \\
\(U_{d,t}\) & 计划购电物理上界 & 防止负价格情景下无意义超量购电 \\
\(z_{d,t}\) & 调增/调减互斥二元变量 & 负价格下连续松弛不能自动保证单边调整 \\
\bottomrule
\end{tabular}
\caption{问题三新增的多阶段调整变量}
\label{tab:6}
\end{table}

物理上界构造为

\[
U_{d,t}
=
\max_{\omega}\overline L_{d,t}^{(\omega)}
+
P_{\max},
\]

即正常计划购电最多用于补足最大情景负荷缺口并以最大功率给储能充电。

问题四中利用

\[
0\le\Delta^+\le U z,
\qquad
0\le\Delta^-\le U(1-z),
\qquad
z\in\{0,1\}
\]

严格保证调增与调减互斥。

\subsection{符号继承关系}

全篇变量不是每问重新定义，而是逐层继承：

\[
\boxed{
\begin{array}{l}
Q1:\ G^L,G^{ch},PV^{ch},C,D,E,V,u\\
Q2:\ Q1 + G^{plan},H,W\\
Q3:\ Q2 + P,A^s,B,\Delta^+,\Delta^-\\
Q4:\ Q3 + \hat\pi,\pi^{(\omega,s)},U,z
\end{array}
}
\]

其中 Q3 中 \(P,A^s,B\) 是对“正常计划购电变量”的阶段化表示，并不与 Q2 的物理执行变量重复。
